\documentclass[11pt]{article}
\usepackage[a4paper,margin=27mm]{geometry}
\usepackage[T1]{fontenc}
\usepackage{lmodern}
\usepackage{microtype}
\usepackage{amsmath,amssymb,amsthm,mathtools}
\usepackage{booktabs}
\usepackage{enumitem}
\usepackage{xcolor}
\usepackage[colorlinks=true,linkcolor=blue!55!black,citecolor=blue!55!black,urlcolor=blue!55!black]{hyperref}
\hypersetup{
  pdftitle={Rank Bounds and a Gauss-Fiber Alternative for Point-Span Tangent Absorption},
  pdfauthor={Lluis Eriksson},
  pdfsubject={Algebraic geometry; first jets; Gauss maps; tangent spaces},
  pdfkeywords={Gauss map, tangent absorption, Veronese re-embedding, Terracini loci, first jets}
}

\newtheorem{theorem}{Theorem}[section]
\newtheorem{proposition}[theorem]{Proposition}
\newtheorem{lemma}[theorem]{Lemma}
\newtheorem{corollary}[theorem]{Corollary}
\newtheorem{remark}[theorem]{Remark}
\newtheorem{definition}[theorem]{Definition}
\newcommand{\PP}{\mathbf P}
\newcommand{\CC}{\mathbf C}
\newcommand{\cO}{\mathcal O}
\newcommand{\Span}{\operatorname{Span}}
\newcommand{\Sym}{\operatorname{Sym}}
\newcommand{\Gr}{\operatorname{Gr}}

\title{Rank Bounds and a Gauss-Fiber Alternative\\
for Point-Span Tangent Absorption}
\author{Lluis Eriksson\\\small Independent researcher}
\date{August 24, 2026\quad (version 0.5)}

\begin{document}
\maketitle

\begin{abstract}
Let $X^d\subset\PP(V)$ be a smooth complex projective variety, let
$H=\cO_X(1)$, and embed $X$ by the complete system $H^m$, where $m\geq3$.
A nonempty finite reduced set $Z\subset X$ is point-span tangent-absorbing
if the vector span $S_Z$ of its evaluation functionals contains the affine
embedded tangent space at every point of $Z$.  We prove
\[
 \dim_{\CC}S_Z,\ |Z|\ \geq
 \min\left\{(d+1)(d+2),\binom{d+m}{d}\right\}.
\]
The proof gives two nonexclusive alternatives.  Either $Z$ contains $d+2$
projectively independent supports whose double points impose independent
conditions, or $Z$ lies in a common projective tangent $d$-plane and in one
fiber of the original Gauss map.  In the second branch, a characteristic-zero
minimal-degree argument shows that the points are unisolvent in degree $m$ on
that plane.  Consequently the stronger bound $(d+1)(d+2)$ is unconditional
for $d\geq3$, and also for $d=2$, $m\geq4$.  For $d=2$, $m=3$, the exact
 minimum is $10$, attained by a triangular lattice of ten points in a linearly
 embedded plane.  Three points on a linearly embedded $\PP^1$ for $m=2$ show
 that the degree assumption cannot be weakened in the uniform theorem.  A
 separate consequence of standard jet-ampleness gives a further lower bound
 that grows linearly with $m$ and eventually improves the constant branch.
\end{abstract}

\section{Introduction}

Throughout, a variety is integral and all schemes and vector spaces are over
$\CC$.  Let $X^d\subset\PP(V)$ be smooth and let $H=\cO_X(1)$.  Write
\[
 \varphi_m=\varphi_{|H^m|}:X\longrightarrow\PP\bigl(H^0(X,H^m)^*\bigr)
\]
for the complete embedding defined by $H^m$, and put
\[
 S_Z:=\Span_{\CC}\{\operatorname{ev}_p:p\in Z\}
      \subset H^0(X,H^m)^*.
\]
We study nonempty finite reduced sets $Z\subset X$ satisfying
\begin{equation}\label{eq:absorption-intro}
 \widehat T_{\varphi_m(p)}\varphi_m(X)\subseteq S_Z\qquad(p\in Z),
\end{equation}
where the left side is the affine tangent space to the cone over the
$H^m$-embedding.  We use $\varphi_m$, rather than an ambient Veronese
notation, because restriction of degree-$m$ forms on $\PP(V)$ need not be
the complete linear series when $X$ is not projectively normal.  We call
\eqref{eq:absorption-intro}
\emph{point-span tangent absorption}.

The condition is an extreme interpolation defect.  The annihilator of the
point span is the space of sections vanishing on $Z$, whereas the annihilator
of the tangent span is the space of sections vanishing on the union $2Z$ of
the first infinitesimal neighborhoods.  Thus absorption is equivalent to
\begin{equation}\label{eq:kernel-equality}
 H^0(X,\mathcal I_Z\otimes H^m)
 =H^0(X,\mathcal I_{2Z}\otimes H^m).
\end{equation}

Write
\[
 \gamma_X:X\longrightarrow\Gr(d+1,V),\qquad
 p\longmapsto\widehat T_pX
\]
for the Gauss map of the original embedding, and define
\begin{equation}\label{eq:gauss-multiplicity}
 g_X(Z):=\max_{W\in\Gr(d+1,V)}|Z\cap\gamma_X^{-1}(W)|.
\end{equation}

The main result is the following.

\begin{theorem}[Rank bound and Gauss-fiber alternative]\label{thm:main}
Let $X^d\subset\PP(V)$ be a smooth complex projective variety, $d\geq1$,
and let $m\geq3$.  If $Z\subset X$ is a nonempty finite reduced
point-span tangent-absorbing set for the complete $H^m$-embedding, then
$Z$ contains $d+1$ supports whose first infinitesimal neighborhoods impose
independent conditions on $H^m$.  In particular,
\begin{equation}\label{eq:baseline}
 \dim S_Z\geq(d+1)^2,\qquad |Z|\geq(d+1)^2.
\end{equation}
Moreover, at least one of the following nonexclusive alternatives holds:
\begin{enumerate}[label=(\alph*)]
\item $Z$ contains $d+2$ supports whose first infinitesimal neighborhoods
impose independent conditions on $H^m$; hence
\[
 \dim S_Z,\ |Z|\geq(d+1)(d+2);
\]
\item there is a vector subspace $W\subset V$, $\dim W=d+1$, such that
\[
 Z\subset X\cap\PP(W),\qquad \widehat T_pX=W\quad(p\in Z),
\]
and in this case
\[
 \dim S_Z,\ |Z|\geq\binom{d+m}{d}.
\]
\end{enumerate}
Consequently,
\begin{equation}\label{eq:main-bound}
 \boxed{\quad
 \dim S_Z,\ |Z|\geq
 \min\left\{(d+1)(d+2),\binom{d+m}{d}\right\}.
 \quad}
\end{equation}
\end{theorem}

The two numerical terms come from different mechanisms.  The first is a
finite first-jet independence statement.  We derive the relevant statement
from Ballico--Chiantini's criterion \cite[Theorem~3.5]{BallicoChiantini2021}
and retain a direct squared-hyperplane proof of this specialization.  The
second term is the dimension of degree-$m$ forms on the common tangent
$\PP^d$ in the Gauss branch.  Its proof uses characteristic zero essentially.

\begin{corollary}[Numerical consequences]\label{cor:numerical}
Under the hypotheses of Theorem~\ref{thm:main}:
\begin{enumerate}[label=(\roman*)]
\item if $d\geq3$, then
$\dim S_Z,|Z|\geq(d+1)(d+2)$ for every $m\geq3$;
\item if $d=2$ and $m\geq4$, then
$\dim S_Z,|Z|\geq12$;
\item if $d=2$ and $m=3$, then $\dim S_Z,|Z|\geq10$, and this is exact;
\item for every $d\geq2$ equality in the baseline
$\dim S_Z=(d+1)^2$ is impossible.
\end{enumerate}
\end{corollary}

\begin{corollary}[Gauss-fiber threshold]\label{cor:gauss-threshold}
If
\[
 g_X(Z)<\binom{d+m}{d},
\]
then alternative~\textup{(a)} of Theorem~\ref{thm:main} holds.  In
particular, $\dim S_Z,|Z|\geq(d+1)(d+2)$.
\end{corollary}

\begin{corollary}[Jet-ampleness supplement]\label{cor:jet-supplement}
Under the hypotheses of Theorem~\ref{thm:main}, set
\[
 J(d,m):=\left\lfloor\frac{m+1}{2}\right\rfloor(d+1)
          +\begin{cases}1,&m\text{ even},\\0,&m\text{ odd}.
            \end{cases}
\]
Then
\[
 \dim S_Z,\ |Z|\geq
 \max\left\{J(d,m),
 \min\left((d+1)(d+2),\binom{d+m}{d}\right)\right\}.
\]
In particular, $J(d,m)>(d+1)(d+2)$ for every $m\geq2d+4$.
This is a lower bound, not an assertion of exactness.
\end{corollary}

The paper makes no claim that the framework of first-jet independence is
new.  The contribution asserted here is limited to the stated deduction for
point-span absorption and the elementary estimate in its exceptional Gauss
branch.  The bibliographic comparison in Section~\ref{sec:comparison} is
selective rather than a certification of priority.

\section{First jets and absorption}\label{sec:setup}

For a smooth point $p\in X$, let $2p$ be the first infinitesimal neighborhood
defined by $\mathfrak m_p^2$.  It has length $d+1$.  Since $H^m$ is very
ample, restriction to $2p$ is surjective.  Intrinsically, its dual is the
affine tangent space to the cone over the $H^m$-embedding; this description
avoids any coordinate-dependent scalar in a symmetric-tensor formula.

\begin{definition}
A nonempty finite reduced set $Z\subset X$ is \emph{point-span
tangent-absorbing with respect to $H^m$} if
$\widehat T_{\varphi_m(p)}\varphi_m(X)\subseteq S_Z$ for every $p\in Z$.
\end{definition}

\begin{lemma}[Kernel criterion]\label{lem:kernel}
The following are equivalent:
\begin{enumerate}[label=(\roman*)]
\item $Z$ is point-span tangent-absorbing;
\item equation \eqref{eq:kernel-equality} holds;
\item restriction to $Z$ and restriction to $2Z$ have the same image rank.
\end{enumerate}
\end{lemma}

\begin{proof}
The annihilator of $S_Z$ is the kernel of restriction to $Z$.  The
annihilator of the span of the affine tangent spaces is the kernel of
restriction to $2Z$.  Since the point span is contained in the tangent span,
the three conditions are equivalent.
\end{proof}

\section{Independent supports and first-jet blocks}\label{sec:independence}

We first isolate the precise independence input.  This is a specialization
of a published general criterion, not a new general theorem about jets.

\begin{lemma}[Projectively independent supports]\label{lem:independent-jets}
Let $S\subset X$ be a finite set whose representatives in $V$ are linearly
independent.  If $m\geq3$, then $2S$ imposes independent conditions on
$H^m$.
\end{lemma}

\begin{proof}
We first explain the specialization of Ballico--Chiantini
\cite[Definition~3.1 and Theorem~3.5]{BallicoChiantini2021}.  Their condition
$\dagger$ for $(X,L_1,S)$ says that, for every $p\in S$, the restriction map
\[
 H^0(X,L_1)\longrightarrow H^0(S\cup2p,L_1|_{S\cup2p})
\]
is surjective.  Take
\[
 L_1=H^2,\qquad L_2=H^{m-2},\qquad L_1\otimes L_2=H^m.
\]
For each $q\in S$, projective independence gives a linear form $h_q$ that
vanishes on $S\setminus\{q\}$ and not at $q$.  The sections
$h_q^{m-2}$ show that $S$ imposes independent conditions on $L_2$.

Fix $p\in S$.  For $q\neq p$, the square $h_q^2$ prescribes the value at
$q$ and vanishes on $2p$.  The square $h_p^2$ supplies the value at $p$.
If $a\in V^*$ vanishes at $p$, then $h_pa$ vanishes at every support and
has differential $h_p(p)\,da$ at $p$.  Because $H$ embeds $X$, the
restrictions of such differentials span the cotangent space of $X$ at $p$.
These quadrics therefore give surjectivity on $S\cup2p$, so
$(X,H^2,S)$ satisfies $\dagger$.  Ballico--Chiantini's Theorem~3.5 now
implies that $2S$ imposes independent conditions on $H^m$.

For completeness, the same specialization has the following direct
squared-hyperplane proof.  Order $S=\{p_1,\ldots,p_r\}$.  A single double
point is independent because $H^m$ is very ample.  Suppose the doubles at
$p_1,\ldots,p_{j-1}$ are independent.  Choose a linear form $e$ vanishing
on their representative span and nonzero at $p_j$.  Then $e^2$ vanishes on
all old double points and is a unit on $2p_j$.  Since $H^{m-2}$ is very
ample, multiplication by $e^2$ maps a space of sections surjectively onto
the first jets at $p_j$.  It can therefore correct arbitrary residual data
at $2p_j$ without changing the old data.  Induction proves the assertion.
This direct proof is included to keep the exact specialization
self-contained.
\end{proof}

This lemma concerns the displayed finite set of supports and retains a direct
proof tailored to that configuration.  The separate global fact that $H^m$ is
$m$-jet ample yields Corollary~\ref{cor:jet-supplement}; its source and role are
made explicit in the proof of that corollary.

\section{Escape or Gauss-fiber concentration}\label{sec:escape}

The next separator is formulated with first jets.  It uses only ambient
linear forms and remains valid without projective normality.

\begin{lemma}[Tangent escape separator]\label{lem:separator}
Let $0\neq v\in W\subsetneq V$ represent $p=[v]\in X$, and let
$u\in\widehat T_pX\setminus W$.  For every $m\geq1$ there is a section of
$H^m$ that vanishes at every point of $X\cap\PP(W)$ but has nonzero first
jet at $p$ in the direction $u$.
\end{lemma}

\begin{proof}
Choose $\lambda\in V^*$ with $\lambda|_W=0$ and $\lambda(u)\neq0$, and
choose $\mu\in V^*$ with $\mu(v)\neq0$.  The restriction to $X$ of
$f=\lambda\mu^{m-1}$ vanishes on $X\cap\PP(W)$, while
\[
 df_v(u)=\lambda(u)\mu(v)^{m-1}\neq0.
\]
Thus its first jet separates the indicated tangent direction from every
linear combination of evaluations supported on $X\cap\PP(W)$.
\end{proof}

\begin{lemma}[Support alternative]\label{lem:support-alternative}
Assume $Z$ is point-span tangent-absorbing.  Let
$p_1,\ldots,p_r\in Z$ have linearly independent representatives and put
$W=\Span(v_1,\ldots,v_r)$.
\begin{enumerate}[label=(\roman*)]
\item If $r\leq d$, some point of $Z$ lies outside $\PP(W)$.
\item If $r=d+1$, either some point of $Z$ lies outside $\PP(W)$, or
\[
 Z\subset\PP(W),\qquad \widehat T_pX=W\quad(p\in Z).
\]
\end{enumerate}
\end{lemma}

\begin{proof}
Suppose $Z\subset\PP(W)$.  If $r\leq d$, then
$\dim W=r<d+1=\dim\widehat T_{p_1}X$, so there is
$u\in\widehat T_{p_1}X\setminus W$.  Lemma~\ref{lem:separator} gives a
section annihilating $S_Z$ but not the corresponding tangent functional,
contrary to absorption.

Now let $r=d+1$.  Both $W$ and every $\widehat T_pX$ have dimension $d+1$.
If $\widehat T_pX\neq W$ for some $p\in Z$, choose
$u\in\widehat T_pX\setminus W$ and apply the same separator.  Hence the
no-escape case forces $\widehat T_pX=W$ for every $p\in Z$.
\end{proof}

\section{Unisolvence in the Gauss branch}\label{sec:gauss-unisolvence}

The following lemma is the additional characteristic-zero input.

\begin{lemma}[Gauss-branch unisolvence]\label{lem:gauss-unisolvence}
Let $W$ have dimension $d+1$, let $Z\subset\PP(W)$ be a nonempty finite
reduced set, and suppose
\[
 \widehat T_pX=W\qquad(p\in Z).
\]
If $Z$ is point-span tangent-absorbing for $H^m$, then, for
$R=\Sym(W^*)$ and the homogeneous ideal $I=I(Z)\subset R$,
\[
 I_m=0.
\]
Consequently,
\[
 \dim S_Z\geq\binom{d+m}{d},\qquad
 |Z|\geq\binom{d+m}{d}.
\]
\end{lemma}

\begin{proof}
Let $F\in I_m$.  The inclusion $W\subseteq V$ induces a surjection
$V^*\twoheadrightarrow W^*$ and hence a surjection
$\Sym^m(V^*)\twoheadrightarrow\Sym^m(W^*)$, so $F$ extends to a degree-$m$
form on $V$.  Its restriction to $X$ vanishes on $Z$.  Under the natural
pairing, the resulting section pairs to zero with every evaluation in
$S_Z$.  Absorption implies that its first jet vanishes along $X$ at every
$p\in Z$.  Because $\widehat T_pX=W$, this is the full first jet of $F$ on
$\PP(W)$.

To identify that condition scheme-theoretically, choose homogeneous
coordinates $x_0,\ldots,x_d$ on $W$ with $p=[1:0:\cdots:0]$.  On the chart
$x_0=1$, vanishing of the value and of all projective directional derivatives
says that $F(1,x_1,\ldots,x_d)$ has no terms of degree zero or one.  It
therefore lies in $(x_1,\ldots,x_d)^2$; by homogeneity, this is equivalent to
$F\in I(p)^2$.  Consequently
\begin{equation}\label{eq:double-ideal}
 I_m\subseteq\bigcap_{p\in Z} I(p)^2.
\end{equation}

Suppose $I_m\neq0$.  Choose a nonzero homogeneous $F\in I$ of minimal
positive degree $\alpha$.  Then $\alpha\leq m$.  For each $p\in Z$, choose
$G_p\in R_{m-\alpha}$ with $G_p(p)\neq0$; when $\alpha=m$, take a nonzero
constant.  Since $G_pF\in I_m$, equation \eqref{eq:double-ideal} gives
$G_pF\in I(p)^2$.  The germ of $G_p$ is a unit at $p$, hence
$F\in I(p)^2$ in the local ring at $p$.

It follows that every first partial derivative of $F$ vanishes at every
point of $Z$.  Thus each partial derivative belongs to the radical ideal
$I(Z)$.  These derivatives have degree $\alpha-1$, and in characteristic
zero at least one is nonzero because $F$ has positive degree.  This
contradicts the minimality of $\alpha$.  Hence $I_m=0$.

The evaluation map
\[
 R_m\longrightarrow\CC^Z
\]
is therefore injective and has rank
$\dim R_m=\binom{d+m}{d}$.  Every form on $W$ extends to $V$, so this value
matrix is a submatrix of the restriction map for $H^m$.  Its rank is at
most both $\dim S_Z$ and $|Z|$, proving the inequalities.
\end{proof}

\begin{remark}\label{rem:characteristic}
Characteristic zero is used at the exact point where a positive-degree
polynomial is required to have a nonzero partial derivative.  In
characteristic $p$, a $p$th power defeats that step.  The separator and the
squared-hyperplane construction are scheme-theoretic and appear
characteristic-free, but Theorem~\ref{thm:main} is not asserted here in
positive characteristic.
\end{remark}

\section{Proof of the main theorem}\label{sec:mainproof}

\begin{proof}[Proof of Theorem~\ref{thm:main}]
Start with any point of $Z$.  Repeated application of
Lemma~\ref{lem:support-alternative}(i) selects $d+1$ points with linearly
independent representatives.  Lemma~\ref{lem:independent-jets} shows that
their double points impose independent conditions on $H^m$.

Dualizing the restriction map to these double points embeds a direct sum of
$d+1$ vector spaces of dimension $d+1$ in $H^0(X,H^m)^*$.  Each summand is
an affine tangent space and hence lies in $S_Z$ by absorption.  Therefore
\[
 \dim S_Z\geq(d+1)^2.
\]
Since $S_Z$ is spanned by $|Z|$ evaluation vectors, also
$|Z|\geq\dim S_Z$.  This proves \eqref{eq:baseline}.

Apply Lemma~\ref{lem:support-alternative}(ii) to the selected $d+1$ points.
If another point escapes their representative span, the resulting $d+2$
supports are projectively independent.  Lemma~\ref{lem:independent-jets}
then gives alternative~(a) and the bound $(d+1)(d+2)$.

If no point escapes, Lemma~\ref{lem:support-alternative}(ii) gives the
subspace $W$ and the common Gauss fiber in alternative~(b).  Apply
Lemma~\ref{lem:gauss-unisolvence} to obtain the bound
$\binom{d+m}{d}$.  Since at least one branch holds, taking the smaller of
their two bounds proves \eqref{eq:main-bound}.
\end{proof}

\begin{proof}[Proof of Corollary~\ref{cor:numerical}]
For $d\geq3$ and $m\geq3$,
\[
 \binom{d+m}{d}\geq\binom{d+3}{d}
 =\frac{(d+1)(d+2)(d+3)}6
 \geq(d+1)(d+2).
\]
For $d=2$ and $m\geq4$, one has
$\binom{m+2}{2}\geq15>12$.  For $(d,m)=(2,3)$ the two terms in
\eqref{eq:main-bound} are $12$ and $10$; exactness is proved in
Proposition~\ref{prop:p2-sharpness} below.

Finally, $(d+1)(d+2)>(d+1)^2$, while for $d\geq2$
\[
 \binom{d+m}{d}\geq\binom{d+3}{d}>(d+1)^2,
\]
the last strict inequality being equivalent to $d(d-1)>0$.  Thus equality
in the square baseline is impossible for $d\geq2$.
\end{proof}

\begin{proof}[Proof of Corollary~\ref{cor:gauss-threshold}]
In alternative~(b), all of $Z$ lies in a single Gauss fiber, and
Lemma~\ref{lem:gauss-unisolvence} gives
$g_X(Z)=|Z|\geq\binom{d+m}{d}$.  The displayed strict hypothesis therefore
excludes alternative~(b), so alternative~(a) must hold.
\end{proof}

\begin{proof}[Proof of Corollary~\ref{cor:jet-supplement}]
A very ample line bundle is $1$-jet ample, and
Beltrametti--Di Rocco--Sommese
\cite[Proposition~2.3]{BeltramettiDiRoccoSommese1999} prove that the tensor
product of an $a$-jet ample bundle and a $b$-jet ample bundle is
$(a+b)$-jet ample.  Thus $H^m$ is $m$-jet ample.  In their convention, a
simple point has weight one, a first neighborhood $2p$ has weight two, and
surjectivity is required for every collection of positive weights summing to
$m+1$.

Suppose first that $m=2q+1$.  If $|Z|\leq q$, enlarge $2Z$ at one support to
a fat-point scheme of total weight $m+1$.  Surjectivity onto that larger
scheme implies surjectivity onto $2Z$.  The dual tangent blocks would then be
independent of total dimension $|Z|(d+1)$, but absorption places their span
inside $S_Z$, whose dimension is at most $|Z|$.  This is impossible because
$d\geq1$ and $Z$ is nonempty.  Hence $Z$ contains $q+1$ distinct supports.
Their first neighborhoods have total weight $2(q+1)=m+1$, so their tangent
blocks are independent and contained in $S_Z$.  Therefore
\[
 \dim S_Z,|Z|\geq(q+1)(d+1)=J(d,m).
\]

If $m=2q$, the same thickening argument excludes $|Z|\leq q$.  Choose $q$
supports and one further support.  The union of the first neighborhoods of
the first $q$ points and the simple last point has total weight
$2q+1=m+1$.  Jet ampleness makes the corresponding dual blocks independent;
they lie in $S_Z$ by absorption and by the definition of the point span.
Consequently
\[
 \dim S_Z,|Z|\geq q(d+1)+1=J(d,m).
\]
Combining this estimate with Theorem~\ref{thm:main} gives the displayed
maximum.  Finally, direct substitution in the even and odd cases shows that
$J(d,m)>(d+1)(d+2)$ for every $m\geq2d+4$.
\end{proof}

\section{Exact examples and the degree threshold}\label{sec:examples}

\begin{proposition}[Exact plane-cubic minimum]\label{prop:p2-sharpness}
For a linearly embedded $X=\PP^2$ and $m=3$, the minimum cardinality and
minimum vector-span dimension of a nonempty reduced point-span
tangent-absorbing set are both $10$.  Equality is attained by
\[
 Z=\{[1:i:j]:i,j\in\mathbf Z_{\geq0},\ i+j\leq3\}.
\]
\end{proposition}

\begin{proof}
The lower bound $10$ is Corollary~\ref{cor:numerical}(iii).  The displayed
set has ten points.  To prove that its evaluations span
$H^0(\PP^2,\cO(3))^*$, dehomogenize and let $f(x,y)$ have total degree at
most three and vanish on the triangular lattice.

On the layer $y=0$, the polynomial $f(x,0)$ has the four roots
$0,1,2,3$, hence vanishes identically and $f=yf_1$ with
$\deg f_1\leq2$.  On $y=1$, the polynomial $f_1(x,1)$ has the three roots
$0,1,2$, so $f_1=(y-1)f_2$ with $\deg f_2\leq1$.  On $y=2$,
$f_2(x,2)$ has the two roots $0,1$, so $f_2=(y-2)f_3$ with $f_3$ constant.
The last point on $y=3$ forces $f_3=0$.  Thus $f=0$.

The evaluation map between two ten-dimensional vector spaces is an
isomorphism.  Hence the point span is the full ambient dual and contains
every tangent space.  Therefore $Z$ is absorbing and
$\dim S_Z=|Z|=10$.
\end{proof}

\begin{lemma}[Simplex-lattice unisolvence]\label{lem:simplex-unisolvence}
For integers $n,r\geq0$, the set
\[
 \Delta_{n,r}:=\{(a_1,\ldots,a_n)\in\mathbf Z_{\geq0}^n:
                         a_1+\cdots+a_n\leq r\}
\]
is unisolvent for polynomials in $n$ variables of total degree at most $r$.
\end{lemma}

\begin{proof}
We induct on $n+r$.  The cases $n=0$ and $r=0$ are immediate.  Suppose that
$f$ of degree at most $r$ vanishes on $\Delta_{n,r}$.  Its restriction
$f(x_1,\ldots,x_{n-1},0)$ vanishes on $\Delta_{n-1,r}$ and is zero by
induction.  Hence $f=x_n g$ with $\deg g\leq r-1$.  For every
$(a_1,\ldots,a_{n-1},s)\in\Delta_{n,r}$ with $s\geq1$, one has $g=0$ at
that point.  Thus
\[
 h(x_1,\ldots,x_{n-1},t)
 :=g(x_1,\ldots,x_{n-1},t+1)
\]
vanishes on $\Delta_{n,r-1}$.  Induction gives $h=0$, hence $g=0$ and
$f=0$.  Since $|\Delta_{n,r}|=\binom{n+r}{n}$ equals the dimension of the
polynomial space, the evaluation map is an isomorphism.
\end{proof}

\begin{remark}[A three-dimensional equality fixture]\label{rem:p3-fixture}
For $X=\PP^3$ and $m=3$, the bound is $20$ and is also attained by
\[
 \{[1:i:j:k]:i,j,k\in\mathbf Z_{\geq0},\ i+j+k\leq3\}.
\]
Lemma~\ref{lem:simplex-unisolvence}, with $(n,r)=(3,3)$, proves that these
twenty evaluations are unisolvent for polynomials of total degree at most
three.  Their span is ambient, so the fixture is absorbing.
\end{remark}

\begin{proposition}[Degree two cannot satisfy the uniform baseline]
\label{prop:p1-m2}
Let $X=\PP^1$ in its linear embedding and let $m=2$.  Any three distinct
points form a point-span tangent-absorbing set with
\[
 \dim S_Z=|Z|=3<4=(d+1)^2.
\]
Consequently the hypothesis $m\geq3$ in the uniform baseline cannot be
weakened to $m=2$.
\end{proposition}

\begin{proof}
A binary quadratic vanishing at three distinct points is zero.  Hence
evaluation at the three points is an isomorphism
\[
 H^0(\PP^1,\cO(2))\longrightarrow\CC^3.
\]
The three evaluation vectors span the full ambient dual, so they contain all
embedded tangent spaces and absorption is automatic.  The asserted
dimensions follow.
\end{proof}

For comparison, when $X=\PP^1$ and $m=3$, four distinct points fill
$H^0(\PP^1,\cO(3))^*$ and attain the square baseline $4$.  Thus the square
bound remains exact in dimension one, although Corollary~\ref{cor:numerical}
shows that it is never exact in dimensions at least two.

\section{Smooth quadrics}\label{sec:quadrics}

\begin{corollary}\label{cor:quadrics}
Let $Q^d\subset\PP^{d+1}$ be a smooth complex quadric and let $m\geq3$.
Every nonempty finite reduced point-span tangent-absorbing set for
$\cO_Q(m)$ satisfies
\[
 \dim S_Z,\ |Z|\geq(d+1)(d+2).
\]
\end{corollary}

\begin{proof}
If $B$ is a nondegenerate symmetric bilinear form defining $Q$, then the
affine tangent space at $[v]$ is the polar hyperplane $v^\perp$.  Polarity is
injective, so every Gauss fiber has one point.  Since
$1<\binom{d+m}{d}$, Corollary~\ref{cor:gauss-threshold} applies.
\end{proof}

For a quadric curve $Q^1\cong\PP^1$ one has
$\cO_Q(m)\cong\cO_{\PP^1}(2m)$.  A reduced set is absorbing exactly when
it has at least $2m+1$ points: at that cardinality its evaluations fill the
ambient dual, while with at most $2m$ points one can construct a degree-$2m$
form vanishing simply at any prescribed support.  This sharper
one-dimensional calculation is compatible with, but not implied sharply by,
the dimension-independent quadric bound above.

\section{Relation to adjacent interpolation notions}\label{sec:comparison}

Terracini loci parameterize finite sets for which the union of double points
fails to impose the expected independent conditions
\cite{BallicoChiantini2021,GaluppiEtAl2023}.  Point-span absorption is more
restrictive: the value rank and the double-point rank are equal.  It therefore
produces an extreme Terracini defect whenever the value span is not already
ambient, but ordinary Terracini defect does not imply absorption.

The independence lemma belongs to the established theory of simultaneous
jet separation.  In the terminology of Beltrametti--Di Rocco--Sommese
\cite{BeltramettiDiRoccoSommese1999}, jet ampleness asks for uniform
surjectivity over prescribed zero-dimensional data.  Their
Proposition~2.3 implies that $H^m$ is $m$-jet ample and is exactly the external
input used in Corollary~\ref{cor:jet-supplement}.  By contrast, the proof of
Lemma~\ref{lem:independent-jets} separates only its specific finite union of
first neighborhoods.  Ballico--Chiantini's
condition $\dagger$ and their Theorem~3.5 provide exactly the product
criterion used in Lemma~\ref{lem:independent-jets}.

Base loci and strong base loci of tangential projections concern spans
generated by tangent spaces that contain an additional point or an additional
tangent space \cite{BallicoBrambillaSantarsiero2026}.  There is an exact
translation in one direction.  Since the evaluation at $p$ lies in the affine
tangent space at $\varphi_m(p)$, absorption gives
\[
 S_Z=\Span\{\widehat T_{\varphi_m(p)}\varphi_m(X):p\in Z\}.
\]
If $A\subseteq Z$ is minimal among subsets whose tangent spaces span $S_Z$,
then every $q\in Z\setminus A$ belongs to the strong base locus of $A$ for
$\varphi_m(X)$.  The extra feature here is that this tangent span equals the
span of evaluations of the full support $Z$.  The Veronese classifications in
\cite{BallicoBrambillaSantarsiero2026} therefore give relevant context, but do
not by themselves establish the point-span bounds proved above.  That source
also shows that strong base loci and Terracini loci are not comparable in
general.

Finally, the exceptional branch uses the Gauss map of the original embedding,
not the Gauss map after the $m$th re-embedding.  General fibers of Gauss maps
in characteristic zero and the additional inseparable phenomena in positive
characteristic belong to a broader classical theory; see, for example,
\cite{Furukawa2013,FurukawaIto2017}.  No classification of Gauss-degenerate
varieties or their absorbing configurations is attempted here.

These comparisons delimit the claims.  They do not establish that no related
statement exists elsewhere, and no exhaustive priority assertion is made.

\section{Scope and limitations}\label{sec:scope}

\begin{itemize}
\item Theorem~\ref{thm:main} is stated and proved over $\CC$.  The
minimal-degree derivative step in Lemma~\ref{lem:gauss-unisolvence} is the
specific obstruction to carrying this proof unchanged to positive
characteristic.
\item The theorem bounds rank and cardinality but does not classify absorbing
sets, even inside the common Gauss fiber of alternative~\textup{(b)}.
\item The plane and three-space equality examples fill their ambient value
spaces.  No assertion is made that all equality cases are ambient or that the
examples classify equality.
\item The comparison with the cited literature is deliberately conservative.
It supplies attribution and logical distinctions, not a comprehensive
bibliographic or priority review.
\item Finite computational rank checks can verify explicit fixtures but do
not prove the universal theorem.  The proof is the argument in
Sections~\ref{sec:independence}--\ref{sec:mainproof}.
\end{itemize}

\section{Conclusion}

Point-span tangent absorption forces two complementary forms of rigidity.
Support escape produces $d+2$ independent first-jet blocks, while failure of
escape concentrates the set in a common tangent plane and forces degree-$m$
unisolvence there.  Combining the two branches yields the minimum in
\eqref{eq:main-bound}, the unconditional strong range in
Corollary~\ref{cor:numerical}, and the exact plane-cubic value $10$.
Jet-ampleness supplies the complementary growing estimate $J(d,m)$ of
Corollary~\ref{cor:jet-supplement}.  The argument also isolates why
characteristic zero and the threshold $m\geq3$ enter the present formulation.

\begin{thebibliography}{99}

\bibitem{BallicoBrambillaSantarsiero2026}
E.~Ballico, M.~C.~Brambilla, and P.~Santarsiero,
\emph{On the strong base locus of a projective variety},
arXiv:2603.15103 (2026).
\url{https://arxiv.org/abs/2603.15103}

\bibitem{BallicoChiantini2021}
E.~Ballico and L.~Chiantini,
\emph{On the Terracini locus of projective varieties},
Milan J. Math. \textbf{89} (2021), 1--17.
\url{https://doi.org/10.1007/s00032-020-00324-5}

\bibitem{BeltramettiDiRoccoSommese1999}
M.~C.~Beltrametti, S.~Di Rocco, and A.~J.~Sommese,
\emph{On generation of jets for vector bundles},
Rev. Mat. Complut. \textbf{12} (1999), 27--45.
\url{https://doi.org/10.5209/rev_REMA.1999.v12.n1.17182}

\bibitem{Furukawa2013}
K.~Furukawa,
\emph{On general fibers of Gauss maps in positive characteristic},
arXiv:1310.5387 (2013).
\url{https://arxiv.org/abs/1310.5387}

\bibitem{FurukawaIto2017}
K.~Furukawa and A.~Ito,
\emph{On separable higher Gauss maps},
arXiv:1702.06010 (2017).
\url{https://arxiv.org/abs/1702.06010}

\bibitem{GaluppiEtAl2023}
F.~Galuppi, P.~Santarsiero, D.~A.~Torrance, and E.~T.~Turatti,
\emph{Geometry of first nonempty Terracini loci},
arXiv:2311.09067 (2023).
\url{https://arxiv.org/abs/2311.09067}

\end{thebibliography}

\end{document}
